
\section{The rank condition: a hypothesis test}
For distinct $a,b$, take $R=M_a$, $B=M_b$, on their irreducible core. Then
\[
 \Delta=\tfrac12e_1(e_a^{\mathsf T}-e_b^{\mathsf T}),
 \quad \pi\Delta=\tfrac{\pi_1}{2}(e_a^{\mathsf T}-e_b^{\mathsf T}).
\]
Since $\pi_1>0$, the row span generated by $\pi\Delta$ already equals
the one-dimensional row space of $\Delta$. This is Chapter 5's rank
condition. When $\gcd(a,b)=1$, the core is aperiodic and the synchronization
criterion applies. When the gcd exceeds one, the same rank condition holds
but synchronization is impossible. Thus aperiodicity is indispensable in
the surrounding criterion; a theorem statement separated from its standing
hypotheses would be misleading.
The example $a=2,b=4$ should be used as a routine hypothesis check.
